\section{Conclusiones}

% ── Frame: Conexión con la teoría ────────────────────────────────────────────
% EDITAR: borrador propuesto. Citas tomadas de Articulo1/secciones/3-Conceptos.tex
%         (traducción propia; el original está en los % ORIG: de la guía).
\begin{frame}{Conexión con la teoría}
  % \begingroup en vez de {…}: un grupo tras el título Beamer lo toma como subtítulo.
  \begingroup\footnotesize
  \renewcommand{\arraystretch}{1.1}
  % Sin >{...}: babel-spanish vuelve activos < y >. Se usa \raggedright
  % en cada celda y \tabularnewline para cerrar renglones.
  \begin{tabularx}{\textwidth}{@{}p{0.30\textwidth}X@{}}
    \toprule
    \textbf{Lo que vivimos → concepto} & \textbf{Lo que dice la investigación} \tabularnewline
    \midrule
    \raggedright Dinámica 1: sostener la mirada y la conversación
    → \textbf{Masking}
      & \raggedright\enquote{El camuflaje fue descrito con frecuencia como algo mental, física y emocionalmente agotador} \citep{hull2017} \tabularnewline
    \raggedright Dinámica 2: un solo objeto nítido y el resto difuso
    → \textbf{Hiperfoco}
      & \raggedright\enquote{Durante un estado de hiperfoco hay una percepción disminuida de los estímulos no relevantes para la tarea} \citep{ashinoff2021} \tabularnewline
    \raggedright Dinámica 3: lectura con audio hasta la fase de crisis
    → \textbf{Meltdown / shutdown}
      & \raggedright Detonantes: \enquote{demandas sociales, frustración, vergüenza, dificultades de comunicación, detonantes emocionales y estímulos sensoriales aversivos abrumadores} \citep{phung2021} \tabularnewline
    \bottomrule
  \end{tabularx}
  \par\endgroup

  \begin{block}{Lo que une a las tres dinámicas}
    {\footnotesize
    El burnout autista surge de \enquote{los estresores de la vida (subtemas:
    enmascaramiento, expectativas, manejo de la discapacidad y cambios vitales)
    que conducen a una carga acumulada} \citep{raymaker2020}: la máscara de la
    Dinámica 1 también alimenta la crisis de la Dinámica 3.\par}
  \end{block}
\end{frame}

% ── Frame: Conclusiones ───────────────────────────────────────────────────────
\begin{frame}{Conclusiones: ¿cómo ayudo?}
  \begin{columns}[T]
    \begin{column}{0.48\textwidth}
      \begin{block}{Lo que ayuda}
        \footnotesize
        \begin{itemize}
          \item Entender que la normalidad no es una ley, sino un estereotipo no universal.
          \item No todas las discapacidades son visibles, ni todas las personas con discapacidad son dependientes.
          \item El término \enquote{inclusión} no solo se aplica a las personas con sexualidad diversa, sino a todas las que no encajan en la norma.
          \item La sinceridad y la claridad ante las diferencias pueden evitar la discriminación.
        \end{itemize}
      \end{block}
    \end{column}
    \begin{column}{0.48\textwidth}
      \begin{alertblock}{Lo que no ayuda}
        \footnotesize
        \begin{itemize}
          \item Usar frases como \enquote{Todos somos un poco autistas} o \enquote{Todos somos un poco discapacitados} para justificar la discriminación.
          \item Propiciar etiquetas o estereotipos que excluyan a las personas con discapacidad de actividades sociales, educativas o laborales.
          \item Creer que adaptar entornos es privilegiar a las personas con discapacidad.
          \item Los discursos políticos sin acciones.
        \end{itemize}
      \end{alertblock}
    \end{column}
  \end{columns}
\end{frame}

% ── Frame: Experiencias personales ───────────────────────────────────────────
\begin{frame}{Experiencias personales}
  \begin{columns}[T,onlytextwidth]
    \begin{column}{0.38\textwidth}
      \begin{block}{Médico}
        {\footnotesize\raggedright
        El reconocimiento médico en México es difícil, tortuoso y burocrático.
        Muchas personas con discapacidad no tienen acceso a un diagnóstico
        oficial y, por lo tanto, no pueden acceder a apoyos o servicios.
        Además, el personal médico no está capacitado para identificar a
        personas con trastorno del espectro autista, y muchas veces se les
        etiqueta como personas con problemas de conducta o de aprendizaje.\par}
      \end{block}
    \end{column}
    \begin{column}{0.29\textwidth}
      \begin{block}{Social}
        {\footnotesize\raggedright
        La libertad de creencia da a los demás la posibilidad de aceptar o no
        a las personas con discapacidad. En la mayoría de los casos, la gente
        se aleja cuando te percibe diferente: se te cataloga como una persona
        rara y se te excluye de actividades sociales.\par}
      \end{block}
    \end{column}
    \begin{column}{0.29\textwidth}
      \begin{block}{Laboral-Educativo}
        {\footnotesize\raggedright
        Un diagnóstico médico no garantiza la inclusión laboral o educativa.
        La discriminación es un problema estructural que se manifiesta en la
        falta de oportunidades, la exclusión de actividades y la
        estigmatización de las personas con discapacidad.\par}
      \end{block}
    \end{column}
  \end{columns}
\end{frame}


% ── Frame: Cierre ────────────────────────────────────────────────────────────
\begin{frame}[plain]
  \begin{center}
    \vfill
    {\color{ColorAcento}\rule{0.6\textwidth}{1.5pt}}\\[0.5cm]
    {\Large\bfseries\color{ColorPrincipal} Gracias}\\[0.3cm]
    {\normalsize \Autor}\\[0.1cm]
    {\small \MateriaCorta}\\[0.1cm]
    {\small \Semestre}\\[0.5cm]
    {\color{ColorAcento}\rule{0.6\textwidth}{1.5pt}}
    \vfill
  \end{center}
\end{frame}
